\subsection*{Learning curves for glyphs and bytes}
We trained a two-layer encoder and two-layer autoregressive decoder from scratch on the same 8,000 dialogue batches as the visual writer. Its 4,096-symbol byte-level BPE vocabulary was fitted on training stories and exchanges. Embedding width was 96, with four attention heads, 256-unit feed-forward layers and a tied output projection. The encoder received all three text fields, segment embeddings and sinusoidal positions. Rationale supervision targeted the first and last source subwords of the annotated word interval. The decoder generated greedily up to 96 subwords with explicit termination. Output comparisons used the same held turns, swap assignments and CoQA normalization. We used 5,000 paired whole-story bootstrap samples for differences in exact-answer rates. These intervals characterize variability over the observed stories with model weights fixed.

A second comparison held the architecture, initialization seed, data order, optimization, spacing and word actions fixed. No learned weights or visual word dictionary were inherited. Each letter supplied either a normalized 256-component observation obtained by pooling a $32\times32$ bitmap, or a 256-dimensional one-hot byte observation. An identical seeded projection summarized the letters. Positional signs bound interior letters to their positions before pooling. Four boundary observations and log length completed a 115-component word vector. Both systems trained the same recurrent pointer and glyph/byte speller. We saved weights at updates 100, 800, 2,400 and 8,000, counted energy through each checkpoint, and measured first-reference accuracy on both held splits. Cost to a target is the first measured cumulative training and feature-preparation energy at which both splits attain that target. This procedure bounds the crossing time only at the sampled checkpoints.

For finite-sample loss, an injective render map with a left inverse permits every text predictor to be transported to the rendered inputs at identical risk, and every image predictor to be pulled back to text. If distinct required answers collide under rendering, a deterministic reader must err on at least one. These Lean results give the conditions under which different training outcomes can be attributed to inductive bias or computational cost rather than a change in accessible information. They require no assumption about how efficiently either predictor is learned.

A second theorem gives an exact experimental criterion. For a fixed pictured context let $g(h)$ be the next finished word image or stop under the deterministic system with previous emitted words $h$. The system writes $w_1,\ldots,w_m$ and stops if and only if $g(w_{<t})=w_t$ at every reference position and $g(w_{1:m})$ is stop. The forward implication follows the generated trace; the reverse follows by induction on the remaining answer. A wrong internal route may still yield the correct visible word. We verified the criterion using actual route decisions, selected physical words, spelling completion and stop, counting the first incorrect physical decision. For style changes a separate Lean lemma gives a sufficient condition: if each action score changes by at most $\epsilon$, a clean winner with margin exceeding $2\epsilon$ remains preferred. We did not measure that bound on handwriting.
